\documentclass[9pt]{article}
\usepackage[a4paper,margin=0.25in]{geometry}
\usepackage{amsmath,amssymb,amsfonts,mathtools}
\setlength{\jot}{4pt} 
\usepackage{braket}
\usepackage{slashed}
\usepackage{cancel}


\usepackage{graphicx}
\usepackage{tikz}
\usepackage{pgfplots}


\usepackage{xcolor}
\usepackage{hyperref}
\usepackage{comment}


\newcommand{\CGCket}{|{\rm CGC}\rangle}
\newcommand{\CGCbra}{\langle {\rm CGC}|}

\def\wt#1{\widetilde{#1}}
\def\mf#1{\mathbb{#1}}

\newcommand{\A}[3]{\mf{A}^{#2}_{#3}(#1)}

\numberwithin{equation}{section}

\title{Do the flagged terms cancel in the total? ($\Lambda\to0$, $\vee\to\infty$)}
\author{}
\date{}

\begin{document}
\maketitle
\noindent This note answers the question left open in \texttt{Other\_terms\_limits}: whether the $\log^{2}(k^{+}/\Lambda)$ pieces and the large-distance tails named in the flags cancel in the total cross section. Conventions as there (phase $e^{-ik\cdot(w-w')}$; non-commuting $\rho$, symmetrised two-$\rho$ sector; projectile not assumed colour neutral; no $1/\epsilon$ written). Statements are marked analytic or numerical in the text; the list of what is not established is at the end of each section.

\section{Soft sector: $\Lambda\to0$ and the double logarithms}
\noindent Notation: $L\equiv\log\frac{k^{+}}{\Lambda}$; $c_{0}\equiv\frac{ig^{4}f^{abd}}{8\pi^{5}k^{+}}$; $K(a,b,z)\equiv\frac{(a-z)\cdot(b-z)}{(a-z)^{2}(b-z)^{2}}$, $\mathcal{K}^{i}(r)=r^{i}/r^{2}$; $\mathcal{C}_{\rm III}$, $\mathcal{L}^{a}_{x'w'}$ as in the limits file. ``Tail'' of a term: the coefficient $T$ in (kernel)$\,\to T/z^{2}$ for the soft gluon at a large transverse position $z$ (angular average), as an operator in the remaining coordinates. Status marks: [A] analytic, [N] numerical.

\subsection{Result}
\begin{enumerate}
\item \textbf{No double logarithm survives in the sum.} Every flagged term (I-3, I-4, I-7, III-1, III-3) is, row by row, of the form $\int_{\Lambda}\frac{dp^{+}}{p^{+}}\int_{z}\frac{T_{\rm row}}{z^{2}}\,F_{\rm row}\big(\tfrac{p^{+}}{k^{+}}z\big)$ with a cut-off function $F$ that is $1$ for $\frac{p^{+}}{k^{+}}|z|\ll$ (a fixed transverse scale) and is produced by the \emph{recoil} of the measured gluon (shift of its position by $\frac{p^{+}}{k^{+}}(z-w)$, equivalently the phase $e^{\mp ik\cdot(z-w)p^{+}/k^{+}}$). Each row alone therefore has $\pi L^{2}\,T_{\rm row}+{\cal O}(L)$. The sum of the tails over all recoil rows vanishes identically,
\begin{align}
\Delta\equiv\sum_{\rm recoil\ rows}T_{\rm row}=0\,,
\end{align}
for a general projectile, in the symmetrised two-$\rho$ sector and also in the symmetrised three-$\rho$ sector [N: $SU(2)$, $SU(3)$, $|\Delta|<10^{-7}$ against single-row tails of order $1$--$6$; changing any one weight gives $|\Delta|={\cal O}(1)$]. Hence the coefficient of $L^{2}$ is zero for a general projectile, a fortiori in the colour-singlet projection.
\item \textbf{There is no cut-off at the energy scale} $z^{2}\sim\frac{k^{+}}{p^{+}}r^{2}$: wherever a single term is cut off by a denominator $p^{+}(x-z)^{2}+k^{+}(x-w)^{2}$, the accompanying term of the same bracket restores it [A], see Eq.~\eqref{eq:B2id}. The rows built only from $\mf{B}_{2}$ and $\mf{A}^{(1)}$ (Group I Rows II, IV; Group II Row I; Group III Rows I, II; the four-$\rho$ terms) and Group II Row IV have tails that are \emph{not cut at all}: they multiply $\log\frac{\vee-k^{+}}{\Lambda}$ resp. $L$ exactly and contain no $L^{2}$.
\item What is left at large $z$ is the infrared divergence of the JIMWLK kernel only: $L\int\frac{d^{2}z}{z^{2}}\,C_{A}$ with the tail of the final JIMWLK formula
\begin{align}
C_{A}=\mathcal{K}(x-w)\cdot\mathcal{K}(x'-w')\sum_{\alpha\in\{x',w'\}}\sum_{\beta\in\{x,w\}}\sigma_{\alpha\beta}\Big[-2N_{c}\,\delta^{b'b}\delta^{c'c}+2f^{gb'c'}f^{gbc}\Big]U^{ab'}(\alpha)U^{ab}(\beta)\,\rho^{c'}(x')\rho^{c}(x)\,,
\end{align}
$\sigma_{x'x}=\sigma_{w'w}=+1$, $\sigma_{x'w}=\sigma_{w'x}=-1$ (units $\frac{1}{(2\pi)^{3}}\frac{g^{4}}{16\pi^{5}k^{+}}$) [N: fitted from the final formula, residual $10^{-8}$]. It is minus the Casimir of the pair $\rho^{c'}\rho^{c}$ acting on the LO structure, vanishes for $\langle\rho^{a}\rho^{b}\rangle\propto\delta^{ab}$ [N: trace $<10^{-7}$], and its $U(z)$-dependent part is zero. The recoil rows do not change it ($\Delta=0$).
\end{enumerate}

\subsection{Rows without recoil (not flagged): the tail is not cut}
\noindent For the single-$\rho$ part of $\mf{B}_{2}$ (Group I Row IV, Group II Row I, Group III Rows I, II, III second part) the two terms that carry the $1/z^{2}$ tail add up, at large $z$ and for every $p^{+}$, to [A]
\begin{align}\label{eq:B2id}
\frac{-\hat z^{j}}{|z|}\,\frac{1}{p^{+}(x-z)^{2}+k^{+}(x-w)^{2}}\left\{\frac{1}{p^{+}}\Big[\frac{(x-w)^{i}(z-w)^{j}}{(z-w)^{2}}+\frac{(x-z)^{j}(x-w)^{i}}{2(x-z)^{2}}\Big]-\frac{1}{k^{+}}\frac{(x-z)^{j}(x-w)^{i}}{2(x-w)^{2}}\right\}\to-\frac{\mathcal{K}^{i}(x-w)}{2\,z^{2}\,k^{+}p^{+}}\,,
\end{align}
because $\frac{1}{p^{+}[p^{+}z^{2}+k^{+}r^{2}]}+\frac{z^{2}}{k^{+}r^{2}[p^{+}z^{2}+k^{+}r^{2}]}=\frac{1}{k^{+}p^{+}r^{2}}$. After the $p^{+}$ integral: $\log\big(1+\frac{k^{+}r^{2}}{\Lambda z^{2}}\big)+\log\frac{(\vee-k^{+})z^{2}}{\Lambda z^{2}+k^{+}r^{2}}=\log\frac{\vee-k^{+}}{\Lambda}$. So the left-over terms of these rows have no $1/z^{2}$ tail at all (in agreement with the unflagged verdicts), and the tail sits entirely in the $\log\frac{k^{+}}{\Lambda}$, $\log\frac{\vee}{k^{+}}$ and $\log(1-\frac{k^{+}}{\vee})$ terms (FLAG I-6 is this tail). The same identity holds in Group I Row V second part ($N_{2a}$, $N_{2b}$: denominator $p^{+}(y-x)^{2}+(k^{+}-p^{+})(y-z)^{2}$; the plus-prescription bracket and the last bracket) and in Group III Row III second part (FLAG III-3 is this cancellation): there the tail is cut only by the phase. Group I Row II: $\Lambda$ occurs only in the coefficient $-2L$; nothing to cut. Four-$\rho$ terms: $\log\frac{\vee}{\Lambda}\int_{z}K(u,v,z)$ with no transverse dependence of the logarithm; their symmetrised four- and three-$\rho$ parts cancel, their two-$\rho$ remainder is part of the JIMWLK list and of $C_{A}$; no $L^{2}$. Group II Row IV ($\xi\to1$): in the variables $y,y'$ of the parent gluons the recoil is the same on both sides and drops out of the phase; kernel $K(y,y',z)$, not cut (this is the unwritten pole of the limits file).

\subsection{Recoil rows (flagged)}
\noindent\textbf{Group I Row III, region $z\to w$ (FLAG I-3).} With $\mathcal{G}_{\Lambda}$ of the flag and $\Phi(z)\equiv\frac{(x'-w')\cdot(y-z)}{(x'-w')^{2}(y-z)^{2}}\,\mathcal{C}_{\rm III}(z)$, the substitution $z_{s}=w-\frac{k^{+}}{q^{+}}(z-w)$ gives, up to terms of relative order $\Lambda/k^{+}$ [A],
\begin{align}
-\frac{c_{0}}{2}\int_{z}\frac{\mathcal{G}_{\Lambda}}{(z-w)^{2}}\,\Phi(z)=-c_{0}\int_{\Lambda}^{k^{+}}\frac{dq^{+}}{q^{+}}\int_{z_{s}}K(x,w,z_{s})\;\Phi\Big(w-\frac{q^{+}}{k^{+}}(z_{s}-w)\Big)\,,
\end{align}
since $-\partial_{\log q^{+}}\frac{1}{2}\log\frac{(x-z_{s})^{2}}{(x-w)^{2}}\,\frac{d^{2}z}{(z-w)^{2}}=K(x,w,z_{s})\,d^{2}z_{s}$. Here $q^{+}=k^{+}-p^{+}$ is the momentum of the soft gluon exchanged with the valence charge at $x$, $z_{s}$ its position, and $z$ the position of the hard gluon before the absorption. Check [A]: for constant $\Phi$ and $|z-w|<R$ the right-hand side is $-c_{0}\Phi\int_{\Lambda}\frac{dq^{+}}{q^{+}}2\pi\log\frac{k^{+}R}{q^{+}|x-w|}=-c_{0}\pi\Phi\log^{2}\frac{k^{+}R}{\Lambda|x-w|}$, the exact formula of FLAG I-3. Consequences: (i) the $\log^{2}$ of FLAG I-3 is not a collinear singularity at $z=w$; it is the large-$z_{s}$ tail ($K(x,w,z_{s})\to1/z_{s}^{2}$) cut off by the recoil, $T=-c_{0}\,\Phi(w)$. (ii) The correct split is
\begin{align}
\underbrace{-c_{0}\,L\int_{z}K(x,w,z)\,\frac{(x'-w')\cdot(y-w)}{(x'-w')^{2}(y-w)^{2}}\,\mathcal{C}_{\rm III}\big|_{U(z)\to U(w)}}_{\text{JIMWLK list: the author's form}}\;-c_{0}\int_{0}^{k^{+}}\frac{dq^{+}}{q^{+}}\int_{z}K(x,w,z)\Big[\Phi\Big(w-\frac{q^{+}}{k^{+}}(z-w)\Big)-\Phi(w)\Big]\,.
\end{align}
The author's integrable form ($-\mathcal{K}(x'-w')\cdot\mathcal{K}(y-w)\,K(x,w,z)$ with $U(w)$) is therefore the right $\log\frac{k^{+}}{\Lambda}$ coefficient; the pointwise form $-\mathcal{K}(x'-w')\cdot\mathcal{K}(y-z)/(z-w)^{2}$ with $U(z)$ is the same term written in the recoil variable, where the split into $L\times$(kernel) + left-over does not exist. The difference of the two forms is exactly what is needed. [N: with the author's form the tail of this term is exactly minus the tail of Group III Row V.]

\noindent\textbf{Group I Row III, large $z$ (FLAG I-4).} The whole row (all terms of the checked row with $\Lambda$ kept, $\vee/k^{+}=50$), $z$-independent part of $\mathcal{C}_{\rm III}$: the angular average of $z^{2}\times$kernel is $c_{0}\,\mathcal{K}(x'-w')\cdot\mathcal{K}(x-w)\,S$ with $\partial S/\partial L=+\frac12$ for $\Lambda|z|<k^{+}|x-w|$ and $-\frac12$ for $\Lambda|z|>k^{+}|x-w|$ [N: slopes $0.5000$ and $-0.51\dots-0.56$ at $|z|/|x-w|=35,\,350$, $\Lambda/k^{+}=10^{-1}\dots10^{-9}$]; nothing happens at $\Lambda z^{2}\sim k^{+}(x-w)^{2}$. Analytically the third term of $N_{3b}$ is
$c_{0}\int_{\Lambda}\frac{dp^{+}}{p^{+}}K(y,w,z)\,\mathcal{K}(x'-w')\cdot\mathcal{K}(a)\big[1-\frac{a\cdot a'}{a'^{2}}\big]$, $a=x-w$, $a'=a-\frac{p^{+}}{k^{+}}(x-z)$, and after the angular average $\mathcal{K}(a)\frac{a\cdot a'}{a'^{2}}\to\mathcal{K}(a')$: it is [LO kernel] $-$ [LO kernel with $w$ shifted by the recoil], with tail $K(y,w,z)\to-1/z^{2}$ [A]. So this term is a pure recoil left-over (weight $c_{A}-c_{D}=1$ in units $c_{0}$, i.e. twice the tail of the $\frac12K(x,y,z)$ term of the row), and the uncut part of the row is the $+\frac12$.

\noindent\textbf{Group I Row V (FLAG I-7), Group III Rows III, V, VI (FLAGS III-1, III-3).} All $1/p^{+}$ (resp. $1/q^{+}$) terms carry the overall phase $e^{-ik\cdot(z-\dots)p^{+}/k^{+}}$; with Eq.~\eqref{eq:B2id} the denominators do not cut [A]; the tail of the row is the tail of its JIMWLK-list coefficient $f(0)$, cut by the phase only: $\pi L^{2}T_{\rm row}$ per row. (The slope $-\frac{\pi}{2}$ ``cut off by the denominator'' quoted in FLAG I-7 holds for the plus-prescription term alone, not for the row.) [The plateaux of these rows were derived analytically, not scanned numerically.]

\subsection{Sum over the recoil rows}
\noindent Tails computed from the $\log\frac{k^{+}}{\Lambda}$ coefficients of the checked rows (soft position $=z$), Weyl-reduced to the symmetrised two-$\rho$ part at two sites, by degree in $U(z)$ (max norm, $SU(3)$; units $\frac{1}{(2\pi)^{3}}\frac{g^{4}}{16\pi^{5}k^{+}}$; c.c. included) [N]:
\begin{center}
\begin{tabular}{l|c|ccc}
row (weight $c_{A}-c_{D}$) & weight & $U(z)^{0}$ & $U(z)^{1}$ & $U(z)U^{T}(z)$\\\hline
Group III Row III, first part & 1 & 0 & 1.46 & 0.92\\
Group III Row III, second part & 1 & 0 & 0.68 & 3.42\\
Group III Row V & 1 & 3.21 & 0 & 0\\
Group III Row VI & 1 & 0 & 4.25 & 4.34\\
Group I Row V, first part & 1 & 0.92 & 1.46 & 0\\
Group I Row V, $N_{2a}$ & 1 & 1.41 & 1.95 & 0\\
Group I Row V, $N_{2b\,\rm rest}$ & 1 & 4.14 & 1.81 & 0\\
Group I Row III, $\frac12K(x,y,z)$ term (FLAG I-4) & 2 & 2.17 & 2.12 & 0\\
Group I Row III, $z\to w$ term, author's form (FLAG I-3) & 1 & 3.21 & 0 & 0\\\hline
sum $\Delta$ & & $2\cdot10^{-8}$ & $1\cdot10^{-8}$ & $3\cdot10^{-8}$
\end{tabular}
\end{center}
The $z$-independent parts ($U^{0}$ and $UU^{T}$) cancel separately among the terms whose recoil acts on $w$ and among those whose recoil acts on $w'$; e.g. $N_{2a}=\frac12N[x'w]+\frac12N[w'x]-N[w'w]$, $N_{2b\,\rm rest}=2N[x'x]-\frac32N[x'w]-\frac32N[w'x]+N[w'w]$ in the basis $N[\alpha\beta]=N_{c}\mathcal{K}\cdot\mathcal{K}\,U^{ac}(\alpha)U^{ab}(\beta)\rho^{c}(x')\rho^{b}(x)$. The part linear in $U(z)$ cancels only between the two sides (each side: $2.37$).

\subsection{How the left-over terms have to be written}
\noindent (a) None of the flagged plus-prescription terms has a limit $\Lambda\to0$ by itself after the transverse integration, and neither has its partner $L\,f(0)$ in the JIMWLK list (logarithmic divergence at large $z$): each pair behaves like $\pi L^{2}T_{\rm row}$. (b) Because $\Delta=0$, the sum of the subtraction terms $f(0)$ of the nine entries above has no $1/z^{2}$ tail: \emph{the plus-prescription terms of Group I Row V, Group III Rows III, V, VI and the two flagged terms of Group I Row III (the latter in the form with $z$ = position of the soft gluon, i.e. the author's form) have to stand under one $p^{+}$ and one $z$ integral}; in this combination $\Lambda\to0$ can be taken and no $L^{2}$ and no large-distance divergence is left in the left-over terms. (c) Alternatively each term can be made finite separately by keeping the recoil phase in the subtraction, $f(0)\to f(0)\,e^{\mp ik\cdot(z-w)p^{+}/k^{+}}$ (for Group I Row III: $\Phi(w)\to$ the shifted LO kernel); the sum of the modified subtractions differs from $L\sum f(0)$ by a finite term because $\Delta=0$. This finite term was not computed. (d) The infrared divergence $L\int d^{2}z/z^{2}\,C_{A}$ of the JIMWLK formula for a non-singlet pair is untouched; with it go FLAG I-6 and the BFKL-side tails.

\subsection{Not established}
\noindent (1) The plateaux of the phase-cut rows (Group I Row V, Group III Rows III, V, VI) rest on the analytic large-$z$ expansion; only Group I Row III was scanned numerically with $\Lambda$ kept. (2) For the part linear in $U(z)$ the cancellation is between terms with the recoil on $w$ and on $w'$; whether it is uniform for $|z|\gtrsim\frac{k^{+}}{\Lambda|k|}$ depends on $U(z)$ at large distance (it is if $U(z)$ tends to a direction-independent limit); not decided. (3) The finite remainder of the combination (b)/(c) was not evaluated, and its existence was not tested by a full numerical integration. (4) The tails of single rows that carry a $\frac{2}{\epsilon}$ bracket (Group I Row II $N_{1}$, Row V $N_{2b\,\rm UV}$, Group II Row IV) depend on the convention $\int d^{2}p/p^{2}=0$; only the tail of the final formula, $C_{A}$, is used above. (5) FLAG I-5 (term-by-term integrability on the line $(x-z)^{2}=(x-w)^{2}$) was not re-examined; it does not involve $L^{2}$.

\section{Energetic sector: $\vee\to\infty$ and the large-distance tails}
\noindent\textbf{Scope.} Flags II-1, II-2, III-2, I-6, I-2 and the unflagged rows with a $1/z^{2}$ tail at $\vee\to\infty$ (Group II Row I; Group III Rows I, II; Group I Row II). $\Lambda$ plays no role here. Notation: $a\equiv k^{+}/\vee$, $\mathcal K^{i}(r)=r^{i}/r^{2}$, $c_{0}\equiv\frac{1}{(2\pi)^{3}}\frac{g^{4}}{16\pi^{5}}\frac{1}{k^{+}}$; $[\dots]_{\rm row}$ is the colour structure of the row exactly as written in \texttt{Other\_terms\_checked} (with $U(z)$ explicit, $y$ integrated). Status of each statement is given: (A) analytic, (N) numerical. Scripts: \texttt{tails.py}, \texttt{sum1.py}, \texttt{sum2.py}, \texttt{proj.py}, \texttt{chk\_an.py}, \texttt{chk23.py} (lattice evaluation of the LaTeX of the checked files with one site $z$ at distance $R=10^{4},10^{6}$; SU(2) and SU(3), Weyl ordering by the matrix engine).

\subsection{What the $\log(1-\frac{k^{+}}{\vee})$ terms are}
\noindent (A, from the checked group files.) In every row in which a $\log\big(1-\frac{k^{+}}{\vee}\big)$ multiplies a transverse integral, it stands in the bracket
\begin{align}
\Big[\log\frac{k^{+}}{\Lambda}+\log\frac{\vee}{k^{+}}+\log\Big(1-\frac{k^{+}}{\vee}\Big)\Big]=\int_{\Lambda}^{\vee-k^{+}}\frac{dp^{+}}{p^{+}}
\qquad\text{or}\qquad
\Big[\log\frac{k^{+}}{\Lambda}+\log\Big(1-\frac{k^{+}}{\vee}\Big)\Big]\simeq\int_{\Lambda}^{\vee-k^{+}}\frac{k^{+}\,dp^{+}}{p^{+}(p^{+}+k^{+})}\notag
\end{align}
in front of one and the same structure (Group I: l.84, 88, 135; Group II: l.49, 69 and Row IV l.194--196; Group III: l.47, 58, 75, 90, 108, 123 of the checked files). Hence each such left-over term is $\log(1-a)$ times the structure that carries $\log\frac{k^{+}}{\Lambda}$ in the same row (the row's entry in the JIMWLK list): it is the correction to the upper limit of the soft-gluon integral, $\vee-k^{+}$ instead of $\infty$. In Group III Rows III, V, VI it is exactly the tail of the subtraction term of the plus prescription, $f(0)\int_{\vee-k^{+}}^{\infty}\frac{k^{+}dp^{+}}{p^{+}(p^{+}+k^{+})}=-f(0)\log(1-a)$. The large-$z$ behaviour of these terms is therefore that of the row-by-row JIMWLK structures (FLAG I-6 is one instance; the same holds for all rows of the table below), not a new one.

\subsection{Tails row by row}
\noindent At fixed $w,w'$ and $|z|\to\infty$ the left-over integrand of each row is $\frac{1}{z^{2}}\times$(coefficient)$\times$(colour structure) $+\,{\cal O}(\log z^{2}/|z|^{3})$. The kernels are isotropic at this order, since $\frac{(y-z)\cdot(z-w)}{(y-z)^{2}(z-w)^{2}}\to-\frac{1}{z^{2}}$, $\frac{(y-z)\cdot(x-z)}{(y-z)^{2}(x-z)^{2}}\to+\frac{1}{z^{2}}$. Result (coefficient of $\frac{1}{z^{2}}$, in units of $c_{0}\,e^{-ik\cdot(w-w')}$; c.c.\ to be added for Groups I, III):
\begin{center}
\begin{tabular}{|l|l|l|}
\hline
row & part that survives $\vee\to\infty$ & part $\propto\log(1-a)$ \\ \hline
I, Row II & none (no $z$) & see text \\
I, Row III ($\mathcal K^{ij}$) & $0$ \ ($|z|^{-3}$) & $0$ \ (kernel $\delta^{ij}/2(z-w)^{2}$) \\
I, Row IV & $0$ & $-i\,\mathcal K(x'-w')\!\cdot\!\mathcal K(x-w)\,\mathcal C_{\rm IV}$ \\
II, Row I $N_{1a}$ & $0$ & $+\frac{i}{2}\,\mathcal K(x'-w')\!\cdot\!\mathcal K(x-w)\,[\dots]_{1a}$ \\
II, Row I $N_{1b}$ & $0$ & $-\frac{i}{2}\,\mathcal K(x'-w)\!\cdot\!\mathcal K(x-w')\,[\dots]_{1b}$ \\
II, Row I $N_{2a\dots2f}$ (sum) & $0$ \ (single terms: $\log z^{2}/z^{2}$) & $+\frac12\,\mathcal K(x'-w')\!\cdot\!\mathcal K(x-w)\,[\dots]_{2}$ \\
II, Row IV, $\log(1-a)$ term & --- & $+2N_{c}\,\mathcal K(x'-w')\!\cdot\!\mathcal K(x-w)\,[U(w')-U(x')][U(w)-U(x)]\rho\rho$ \\
III, Row I & $0$ & $+i\,\mathcal K(x'-w')\!\cdot\!\mathcal K(x-w)\,[\dots]_{\rm I}$ \\
III, Row II & $0$ & $+i\,\mathcal K(x'-w')\!\cdot\!\mathcal K(x-w)\,[\dots]_{\rm II}$ \\
III, Row III, 1st part & $0$ (phase in $z$) & $-i f^{c'd'a}\,\mathcal K(x'-w')\!\cdot\!\mathcal K(y-w)\,[\dots]_{\rm III.1}$ \\
III, Row III, 2nd part & $0$ & $+\,\mathcal K(x'-w')\!\cdot\!\mathcal K(x-w)\,[\dots]_{\rm III.2}$ \\
III, Row V & $0$ (phase in $z$) & $-2i f^{c'ca}\,\mathcal K(x'-w')\!\cdot\!\mathcal K(y-w)\,[\dots]_{\rm V}$ \\
III, Row VI & $0$ (phase in $z$) & $-2i f^{c'd'a}\,\mathcal K(x'-w')\!\cdot\!\mathcal K(y-w)\,[\dots]_{\rm VI}$ \\ \hline
\end{tabular}
\end{center}
(in Rows III, V, VI of Group III $y'$ is renamed $w'$ in the $\log(1-a)$ term, which has no $p^{+}$ phase). Status: second column (N): $R^{2}\times$row at $\vee=\infty$ decreases like $\log R/R$ for Group I Row IV, Group II Row I ($N_{1a}$, $N_{1b}$, and the \emph{sum} of $N_{2a},\dots,N_{2f}$; $N_{2a}$, $N_{2b}$, $N_{2c}$ separately grow like $\log R$), Group III Rows I, II, in all Weyl components; for Group I Row III the statement is the one of the limits file ($\mathcal K^{ij}\sim|z|^{-3}$), not re-derived here. Third column: (A) from the kernels, and (N): the tail divided by $\log(1-a)$ is independent of $a$ ($a=0.3$, $0.05$) and agrees with the entries for Group I Row IV, Group II $N_{1a}$, Group III Rows I, II to relative $10^{-6}$; the net BFKL terms of Group III Rows III, V, VI (the third terms in the BFKL file) have \emph{no} $1/z^{2}$ tail (N).

\noindent\textbf{Sum of the tails} (N; SU(2) and SU(3), ordered products reduced to symmetrised ones).
\begin{itemize}
\item At $\vee=\infty$ (coefficient $\log(1-a)$ set to zero) none of the rows of the table has a $1/z^{2}$ tail at fixed $w,w'$: there is nothing to cancel between rows. The only left-over terms with a large-distance tail at $\vee=\infty$ are those of Group II Row IV (next subsection), and the $\Lambda$-dependent ones (FLAG I-4, III-1, I-7; not this part).
\item The $\log(1-a)$ tails: Group I Row IV and Group III Row I cancel each other identically (all Weyl components, to $10^{-14}$ relative). The remaining ones do not cancel: for the sum of all rows of the table the symmetrised three-$\rho$ part, the symmetrised two-$\rho$ part, and the colour-singlet projection of the latter ($\langle\rho^{a}\rho^{b}\rangle\propto\delta^{ab}$), also with the neutrality condition, are all non-zero, of the same size as the single rows (e.g. SU(2): two-$\rho$ norm $1.6\times10^{-6}$ against $0.3$--$1.2\times10^{-6}$ for the rows). This is as it must be: the sum is the tail of the partial sum of the row-by-row JIMWLK structures over those rows whose soft integral ends at $\vee-k^{+}$; the rows whose soft integral ends at $k^{+}$ (Group I Row V, part of Rows II, III, the four-$\rho$ terms) are not in it. The Row II term of Group I, $-\log(1-a)\,\mathcal B\,\mathcal S_{\rm LO}$ with the bracket $\mathcal B$, is $-\frac12\times$ the Group II Row IV entry if $\mathcal B$ is read as $-\frac1\pi\int_{z}$; adding it does not produce a cancellation either.
\end{itemize}
\noindent\textbf{Verdict I-6 (and all $\log(1-a)$ terms).} They are $\log(1-a)={\cal O}(k^{+}/\vee)$ times integrals that diverge logarithmically at large $z$ for a general projectile and also in the singlet projection; the divergences do not cancel in the sum of these terms. They vanish for $\vee\to\infty$ only if the limit is taken with the large-distance regulator in place (dimensional regularisation, or a cut-off $|z|<R$: then they are ${\cal O}\big(\frac{k^{+}}{\vee}\log R\big)$); the limits $\vee\to\infty$ and $R\to\infty$ do not commute. They are not a new structure: each is the row's JIMWLK structure, and should be written with it, as $\big[\log\frac{k^{+}}{\Lambda}+\log(1-\frac{k^{+}}{\vee})\big]\times$(structure), not among the left-over terms. \textbf{Verdict I-2:} only the sum $\mathcal K^{ij}$ is defined; in it $\log(1-a)$ enters as $\log\big|\frac{(x-z)^{2}-a(x-w)^{2}}{(1-a)(x-w)^{2}}\big|\,\frac{\delta^{ij}}{2(z-w)^{2}}$, which vanishes at $z=w$ and has no $1/z^{2}$ tail (with the factor $\mathcal K^{j}(y-z)$); no contribution to the sums above.

\subsection{Group II Row IV: remnant, DGLAP term (FLAG II-1, II-2)}
\noindent Write $C\equiv\frac{1}{(2\pi)^{3}}\frac{g^{4}}{16\pi^{4}}\frac{1}{k^{+}}$, $F(y',y)\equiv\int_{x,x'}\mathcal K(x-y)\!\cdot\!\mathcal K(x'-y')\,[U(y')-U(x')]^{ce'}[U(y)-U(x)]^{ce}\rho^{e'}(x')\rho^{e}(x)$, $\tilde F(s)\equiv\int_{Y}F(Y-\tfrac s2,Y+\tfrac s2)$, and $X\equiv\frac{4\pi}{k^{2}}\tilde F(0)$ (the BFKL structure of this row; $2N_{c}C\,X$ is the coefficient of $\log\frac{\vee}{k^{+}}$).

\noindent\textbf{(a) The remnant is not an independent term} (A). Since $\int_{a}^{1}d\xi\,\frac{1-\xi}{\xi}=\log\frac{\vee}{k^{+}}-(1-a)$ exactly, the compensating term of $N_{\rm DGLAP}$ is $-2N_{c}C\big[\log\frac{\vee}{k^{+}}-(1-a)\big]X$, i.e. exactly minus (BFKL term $+$ remnant) of this row. Therefore
\begin{align}
N_{-(1-k^{+}/\vee)}+N_{\rm DGLAP}\Big|_{\text{compensator }\frac{1-\xi}{\xi}}\;=\;N_{\rm DGLAP}\Big|_{\text{compensator }\frac{1}{\xi}}\,,\notag
\end{align}
the same expression with $2N_{c}\frac{1}{\xi}$ instead of $2N_{c}\frac{1-\xi}{\xi}$ in the second term, and with it the BFKL coefficient of the row is $\log\frac{\vee}{k^{+}}$ without remnant. The $+1$ that matches the $-1$ of the remnant is $-\int_{0}^{1}d\xi\,(-1)$ of the compensator; no other left-over term contains a pure number times $X$.

\noindent\textbf{(b) $\xi\to0$ after the $s=y-y'$ integration} (A; N for model functions). With $q=k/\xi$,
\begin{align}
\int d^{2}s\,e^{-iq\cdot s}\Big[2\gamma_{E}+\log\frac{s^{2}\mu^{2}}{4\xi^{2}}\Big]\tilde F(s)=-\frac{4\pi}{q^{2}}\tilde F(0)-\frac{8\pi i\,q^{m}}{q^{4}}\,\partial_{m}\tilde F(0)+{\cal O}(q^{-4})\,,\notag
\end{align}
from $\int d^{2}s\,e^{-iq\cdot s}\log s^{2}=-\frac{4\pi}{q^{2}}$ ($q\neq0$) and its $q$-derivatives; the constant under the logarithm contributes only the Fourier transform of $\tilde F$ at large $q$ (smaller than any power for smooth $\tilde F$). Hence the integrand of the $\xi$ integral is
\begin{align}
-\frac{4\pi}{k^{2}}\Big[\mathcal P_{gg}(\xi)-\frac{2N_{c}}{\xi}\Big]\tilde F(0)\;-\;\mathcal P_{gg}(\xi)\,\xi\,\frac{8\pi i\,k^{m}\partial_{m}\tilde F(0)}{k^{4}}\;+\;{\cal O}(\xi)\,:\notag
\end{align}
the first term is ${\cal O}(1)$ with the compensator $\frac1\xi$ (and ${\cal O}(\xi)+{\rm const}$, the constant being the remnant, with $\frac{1-\xi}{\xi}$); the second is ${\cal O}(1)$, and it is absent in the symmetrised two-$\rho$ part, where $\tilde F(s)=\tilde F(-s)$. The lower limit can be set to $0$. Smoothness needed: $\tilde F(s)-\tilde F(0)={\cal O}(|s|^{\delta})$, $\delta>0$, is sufficient (remainder ${\cal O}(q^{-2-\delta})$). For differentiable $U$ this holds also for point-like charges: at $y'\to y$ the bilinear tends to $[U(y)-U(x')][U(y)-U(x)]$, the kernels $\mathcal K(x-y)[U(y)-U(x)]$ are bounded at $y\to x$, and $\tilde F$ has at most $|s|\log|s|$ cusps (N: for $\tilde F=e^{-s^{2}}$ the first two terms of the expansion are reproduced and the integrand with compensator is $\propto\xi$; for $\tilde F=e^{-s^{2}}(1+|s|)$ the remainder is ${\cal O}(\log q/q^{3})$, i.e. an integrable $\log\xi$). At fixed $s$ the limit does not exist, also with the compensator written under the same $s$ integral (the $\xi$ integral is then a distribution $\propto\partial_{\sigma}[\pi\delta(\sigma)-i{\rm P}\frac1\sigma]$, $\sigma=k\cdot s$). Convergent form:
\begin{align}
N_{\rm DGLAP}+N_{-(1-k^{+}/\vee)}=-C\int_{0}^{1}d\xi\,\Bigg\{\frac{\mathcal P_{gg}(\xi)}{\xi^{2}}\int d^{2}s\,e^{-ik\cdot s/\xi}\Big[2\gamma_{E}+\log\frac{s^{2}\mu^{2}}{4\xi^{2}}\Big]\tilde F(s)+\frac{2N_{c}}{\xi}\,\frac{4\pi}{k^{2}}\,\tilde F(0)\Bigg\}\,,\notag
\end{align}
$s$ integral first.

\noindent\textbf{(c) Large distance (FLAG II-2)} (A, under the stated assumption). At large $Y$, $F(Y,Y)\to\frac{1}{Y^{2}}\,G_{\infty}(Y)$, $G_{\infty}(Y)\equiv\int_{x,x'}[U(Y)-U(x')]^{ce'}[U(Y)-U(x)]^{ce}\rho^{e'}(x')\rho^{e}(x)$, so $\tilde F(0)$ and $X$ diverge logarithmically unless $G_{\infty}\to0$. $G_{\infty}$ does not vanish for a neutral projectile nor in the singlet projection (with $\langle\rho^{e'}(x')\rho^{e}(x)\rangle\propto\delta^{e'e}G(x',x)$, $\int_{x}G=0$, it is $\int_{x,x'}{\rm tr}[U^{\dagger}(x')U(x)]\,G(x',x)$; the other three terms are removed by neutrality). \emph{If} $U$ varies slowly at large distance (so that $F(Y-\frac s2,Y+\frac s2)\simeq F(Y,Y)$ in the tail), the tails of the left-over terms of this row are pure numbers times the tail of $X$:
\begin{align}
\text{DGLAP}+\text{remnant}:&\quad-C\,\frac{4\pi}{k^{2}}\int_{0}^{1}d\xi\Big[\frac{2N_{c}}{\xi}-\mathcal P_{gg}(\xi)\Big]\frac{G_{\infty}}{Y^{2}}=-C\,\frac{4\pi}{k^{2}}\,\frac{\beta_{0}}{2}\,\frac{G_{\infty}}{Y^{2}}\,,\qquad \beta_{0}=\tfrac{11N_{c}}{3}\notag\\
\text{RC}:&\quad-C\,\frac{4\pi}{k^{2}}\,\frac{\beta_{0}}{2}\,\frac{G_{\infty}}{Y^{2}}\,,\qquad
\text{Finite}:\quad+2N_{c}C\,\frac{4\pi}{k^{2}}\,\frac16\,\cos(2\varphi_{kY})\,\frac{G_{\infty}}{Y^{2}}\notag\\
\text{sum}:&\quad 2N_{c}C\,\frac{4\pi}{k^{2}}\Big[-\frac{11}{6}+\frac16\cos(2\varphi_{kY})\Big]\frac{G_{\infty}(Y)}{Y^{2}}\notag
\end{align}
($\int_{0}^{1}d\xi[\xi(1-\xi)-1+\frac{\xi}{(1-\xi)_{+}}]=-\frac{11}{6}$; the $\delta(1-\xi)$ part of $\mathcal P_{gg}$ and the RC term cancel; $\int_{0}^{1}d\xi\,\xi(1-\xi)=\frac16$; $\varphi_{kY}$ is the angle between $k$ and $Y$). So the left-over terms of Group II Row IV keep a logarithmic large-distance divergence $-\frac{11}{6}\times$(tail of the BFKL-type structure $2N_{c}C\,X$), plus an anisotropic piece that integrates to zero only if $G_{\infty}$ does not depend on the direction. Nothing in the other left-over terms was found that cancels it (Group I has $\frac{11}{3}-2\cdot\frac{11}{6}=0$ times the LO structure). Without the slow-variation assumption the tail is $\frac{1}{Y^{2}}$ times a function of $U(Y\pm\frac s2)$ under the $s$ integral and is not a multiple of the tail of $X$. Not verified numerically.

\subsection{Group III Rows III, V, VI: the subtracted terms (FLAG III-2)}
\noindent (A; N for a model function.) With $r=y'-z$ the second term of these rows is
\begin{align}
\propto\int_{0}^{\vee-k^{+}}\frac{dp^{+}}{2\pi}\Big(-\frac{1}{k^{+}}\Big)\,e^{-ik\cdot(w-z)}\int d^{2}r\,e^{i(1+\frac{p^{+}}{k^{+}})k\cdot r}\,\mathcal K^{m}(r)\big[B^{k'}(z+r)-B^{k'}(z)\big]\,,\notag
\end{align}
$B^{k'}(y')=\mathcal K^{k'}(x'-y')[U(y')\cdots-\cdots U(x')]\rho(x')$. Each term alone gives, at large $p^{+}$, $B(z)\,\frac{2\pi i\,k^{m}}{k^{2}(1+p^{+}/k^{+})}$, hence $\pm\log\frac{\vee}{k^{+}}$; in the difference $\mathcal K^{m}(r)[B(z+r)-B(z)]$ is bounded at $r\to0$ (homogeneous of degree $0$) and the $r$ integral is ${\cal O}(k^{+2}/p^{+2})$ (N: $1/p^{+}$ for one term with the coefficient $2\pi B(z)$, faster than $1/p^{+2}$ in the checked component of the difference). The upper limit can be set to $\infty$ with an error ${\cal O}(k^{+}/\vee)$. Convergent form: with $r_{\parallel}=k\cdot r/|k|$, $r_{\perp}$ integrated first,
\begin{align}
\int_{0}^{\infty}dp^{+}\,e^{i\frac{p^{+}}{k^{+}}k\cdot r}=k^{+}\Big[\pi\,\delta(k\cdot r)+i\,{\rm P}\frac{1}{k\cdot r}\Big]\ \ \Big(=k^{+}\big[\pi\delta(s)-i{\rm P}\tfrac1s\big],\ s=k\cdot(z-y')\Big),\notag\\
\text{2nd term}\;\propto\;-\frac{1}{2\pi}\,e^{-ik\cdot(w-z)}\int dr_{\parallel}\,e^{i|k|r_{\parallel}}\Big[\frac{\pi}{|k|}\delta(r_{\parallel})+\frac{i}{|k|}{\rm P}\frac{1}{r_{\parallel}}\Big]\int dr_{\perp}\,\mathcal K^{m}(r)\big[B^{k'}(z+r)-B^{k'}(z)\big]\,.\notag
\end{align}
(The $r_{\perp}$ integral is ${\cal O}(r_{\parallel})$ for $m=\parallel$ and a principal value for $m=\perp$; the two terms of the bracket must stay together.) \emph{Is $\log\vee$ removed completely?} Yes. The term added back is $B(z)\int_{0}^{\vee-k^{+}}\frac{dp^{+}}{p^{+}+k^{+}}=B(z)\log\frac{\vee}{k^{+}}$ exactly (third term of these rows in the BFKL file), with no constant; the plus-prescription rewriting gives $\mp\log\frac{\vee}{k^{+}}$ times the first structure (first and second BFKL terms, which cancel) and $\log(1-a)$ times it (table above). So for Row III first part and Rows V, VI the finite, $\vee$-independent remainder is the subtracted term itself in the form above, nothing else. For Row III second part the kernel is $\frac{p^{+}}{k^{+2}[p^{+}(x-z)^{2}+k^{+}(x-w)^{2}]}$, and the term added back is $\frac{1}{(x-z)^{2}}\log\frac{\vee}{k^{+}}$ (BFKL) plus the finite remainder $-\int_{0}^{\infty}\frac{dp^{+}\,(x-w)^{2}}{(x-z)^{2}(p^{+}+k^{+})[p^{+}(x-z)^{2}+k^{+}(x-w)^{2}]}$, which is the third term of $N_{2}$ in the list of other terms (already there; nothing is missing). Large $z$: these terms keep the phase $e^{-ik\cdot(w-z)}$ and have no $1/z^{2}$ tail without phase.

\subsection{Pure numbers times LO-, BFKL- or JIMWLK-type structures in the left-over terms}
\begin{center}
\begin{tabular}{|l|l|l|}
\hline
term & number & structure \\ \hline
Group I Row II $N_{1}$ & $-\frac{67}{9}+\frac{\pi^{2}}{3}$ & $\mathcal S_{\rm LO}$ (LO type) \\
Group I Row II $+$ Row V ($\mathcal B$ terms) & $\frac{11}{3}-2\cdot\frac{11}{6}=0$ & $\mathcal B\,\mathcal S_{\rm LO}$ \\
Group II Row IV remnant & $-1$ & $2N_{c}C\,X$ (BFKL type) \\
Group II Row IV compensator, $-\int_{0}^{1}d\xi(-1)$ & $+1$ & $2N_{c}C\,X$ \\
Group II Row I $N_{2f}$, $[\frac1\vee-\frac{1}{k^{+}}]$ & $-1$ & $\frac{g^{4}d_{\perp}}{32\pi^{5}}\int\frac{[\dots]_{2}}{(z-w)^{2}(z-w')^{2}}$: not of LO/BFKL/JIMWLK type \\
all $\log(1-a)$ terms & $\log(1-a)\to0$ & JIMWLK structures of the rows \\
Group III & none & (third BFKL terms exact) \\ \hline
\end{tabular}
\end{center}
Sum: BFKL type $-1+1=0$ (the remnant is absorbed, (a) above); LO type $\big(-\frac{67}{9}+\frac{\pi^{2}}{3}\big)\mathcal S_{\rm LO}$; JIMWLK type none at $\vee=\infty$. The power term of $N_{2f}$ has no $x$-dependent kernel, decays like $1/z^{4}$, and cannot be separated from the other two terms of $N_{2f}$ ($z\to w,w'$); it does not combine with the remnant.

\subsection{Conclusions for this sector}
\begin{enumerate}
\item \emph{Term-by-term $\vee\to\infty$.} (i) Write every $\log(1-\frac{k^{+}}{\vee})$ term together with the $\log\frac{k^{+}}{\Lambda}$ term of the same row (same structure); then no left-over term contains $0\times\infty$, and the limit is taken in the coefficient. (ii) Group I Row III: only the combined kernel $\mathcal K^{ij}$. (iii) Group II Row I: $N_{2a},\dots,N_{2f}$ under one $z$ integral. (iv) Group II Row IV: drop the remnant and use the compensator $\frac{2N_{c}}{\xi}$; $s$ integral before $\xi$ integral. (v) Group III Rows III, V, VI: the subtracted terms in the distributional form, $y'-z$ integral before $p^{+}$ (or $r_{\perp}$ first).
\item \emph{What does not cancel.} The tails of the $\log(1-a)$ terms do not cancel among themselves, in any channel; they vanish with the coefficient only at fixed large-distance regulator. At $\vee=\infty$ the only left-over large-distance divergence found in this sector is that of Group II Row IV, $\propto\big[-\frac{11}{6}+\frac16\cos2\varphi\big]G_{\infty}/Y^{2}$ (under the slow-variation assumption), present also for a neutral, singlet projectile because $G_{\infty}\neq0$ there.
\item \emph{Not established.} Whether the Row IV tail is cancelled by, or has to be attributed to, the large-distance term of the BFKL formula or the unwritten collinear pole (no matching term was found among the left-over terms); the tails in the regions where $w$, $w'$ go to infinity together with $z$ were estimated by power counting only (they are absent for the rows of the table, present for Row IV where they are the $Y$ tail above); the three-$\rho$ rows were evaluated on lattices with two charge sites, configurations $x=x'$, $w=w'$ excluded; the $\frac16\cos2\varphi$ term and the $-\frac{11}{6}$ are analytic, not checked numerically; Group I Row III and the $\mathcal A_{0}$ terms were not re-evaluated.
\end{enumerate}

\end{document}
